\documentclass[12pt]{article}
\def\activityid{F08-M-v3}\usepackage[letterpaper,margin=.65in,top=.60in,bottom=.65in]{geometry}
\usepackage[T1]{fontenc}
\usepackage[default]{sourcesanspro}
\usepackage{microtype,tikz,array,tabularx,fancyhdr,amsmath,amssymb}
\usepackage[hidelinks]{hyperref}
\usetikzlibrary{calc,arrows.meta}
\setlength{\parindent}{0pt}\setlength{\parskip}{7pt}
\setlength{\headheight}{0pt}\setlength{\footskip}{26pt}\setlength{\emergencystretch}{2em}
\pagestyle{fancy}\fancyhf{}\renewcommand{\headrulewidth}{0pt}\renewcommand{\footrulewidth}{.4pt}
\fancyfoot[L]{\scriptsize Bellingham Math Circle / Week 8 / \activityid}
\fancyfoot[R]{\scriptsize \thepage}
\definecolor{ink}{gray}{.12}\definecolor{soft}{gray}{.95}\color{ink}
\newcommand{\PageTitle}[2]{{\fontsize{10}{12}\selectfont\bfseries WEEK 8\quad STRATEGY LAB\hfill #1\par}\vspace{3pt}{\fontsize{26}{29}\selectfont\bfseries #2\par}\vspace{2pt}}
\newcommand{\Subhead}[1]{\par\vspace{3pt}{\large\bfseries #1}\par}
\newcommand{\RuleBox}[1]{\par\begingroup\setlength{\fboxsep}{8pt}\colorbox{soft}{\parbox{\dimexpr\linewidth-16pt\relax}{#1}}\endgroup\par}
\newcommand{\WriteBox}[2]{\par\textbf{#1}\par\vspace{2pt}\begin{tikzpicture}\draw[black!40,line width=.5pt] (0,0) rectangle (\dimexpr\linewidth-.5pt\relax,#2);\end{tikzpicture}\par}
\newcommand{\RookBoard}[3]{\begin{tikzpicture}[x=#3,y=#3]\draw[step=1,black!45] (0,0) grid (#1,#2);\draw[thick] (0,0) rectangle (#1,#2);\node[font=\Large] at ({#1-.5},.5){$\star$};\draw[-{Latex},thick] (0,-.35)--(#1,-.35);\draw[-{Latex},thick] ({#1+.35},#2)--({#1+.35},0);\end{tikzpicture}}
\newcommand{\Table}[3]{\begin{tikzpicture}[x=#3,y=#3]\draw[step=1,black!25] (0,0) grid (#1,#2);\draw[thick] (0,0) rectangle (#1,#2);\fill (0,0)circle(3pt);\draw[-{Latex},very thick] (0,0)--(.7,.7);\node[below] at ({#1/2},-.1){width #1};\node[rotate=90,above] at (-.2,{#2/2}){height #2};\end{tikzpicture}}
\newcommand{\GraphNodes}{\coordinate(A)at(0,0);\coordinate(B)at(3,0);\coordinate(C)at(1.5,2.6);\coordinate(D)at(1.5,.95);}
\newcommand{\Kfour}[1]{\begin{tikzpicture}[scale=#1]\GraphNodes\draw[very thick](A)--(B)--(C)--(A)--(D)--(B) (D)--(C);\foreach \n in {A,B,C,D}{\node[circle,draw,fill=white,minimum size=22pt,font=\bfseries]at(\n){\n};}\end{tikzpicture}}
\newcommand{\House}[1]{\begin{tikzpicture}[scale=#1]\coordinate(A)at(0,0);\coordinate(B)at(3,0);\coordinate(C)at(3,2);\coordinate(D)at(0,2);\coordinate(E)at(1.5,3.3);\draw[very thick](A)--(B)--(C)--(D)--(A) (D)--(E)--(C);\foreach \n in {A,B,C,D,E}{\node[circle,draw,fill=white,minimum size=22pt,font=\bfseries]at(\n){\n};}\end{tikzpicture}}




\newcommand{\StudentPage}[1]{{\fontsize{10}{12}\selectfont\bfseries Week 8 / Strategy lab / #1\par}\vspace{10pt}}
\newcommand{\Problem}[2]{\par\noindent\textbf{Problem #1:} #2\par}
\newcommand{\Space}[1]{\par\begin{tikzpicture}\draw[black!35] (0,0) rectangle (\dimexpr\linewidth-.5pt\relax,#1);\end{tikzpicture}\par}
\newcommand{\Piles}[1]{\begin{center}\begin{tikzpicture}[x=1in,y=1in]\foreach \x in {0,...,#1}{\draw[black!50,rounded corners] ({\x*2.05},0) rectangle ({\x*2.05+1.8},1.3);}\end{tikzpicture}\end{center}}

\begin{document}

\StudentPage{Grades 2--3}
\Problem{1}{Share a token. Take turns moving it right OR down one or more squares. Do not move diagonally. Whoever reaches the star wins. Play the starts below, then swap who starts. A start written (2,1) is 2 squares left and 1 square above the star.}
\begin{center}\RookBoard{4}{4}{.59in}\end{center}
\begin{center}\renewcommand{\arraystretch}{1.8}\begin{tabular}{c|p{2.25in}|p{2.25in}}Start&First game: who won?&After swapping: who won?\\\hline(1,0)&&\\(1,1)&&\\(2,1)&&\\(2,2)&&\end{tabular}\end{center}
\Problem{2}{Replay one start. Choose a different first move. Record that move and the reply.}
\Space{.7in}
\newpage
\StudentPage{Grades 2--3}
\Problem{3}{Play from (3,1), (3,2), (2,3), and (3,3) on the board below. For each, write a first move you would choose. Test it against your partner.}
\begin{center}\renewcommand{\arraystretch}{2}\begin{tabular}{c|p{2.25in}|p{2.25in}}Start&My move leaves&Partner's reply leaves\\\hline(3,1)&&\\(3,2)&&\\(2,3)&&\\(3,3)&&\end{tabular}\end{center}
\Problem{4}{Make a map of the board. Mark T on a square if its next player cannot win against perfect replies. Mark W if that player can force a win, and ? while unsure. Mark the star T: the game has ended, so its next player has no move. Work outward from the star. A W needs one move to a T; a T has only moves to W squares.}
\begin{center}\RookBoard{4}{4}{.63in}\end{center}
\newpage
\StudentPage{Grades 2--3}
\Problem{5}{Put the token at (3,2) on your board. Build piles of 3 and 2 in the spaces below. Slide right 1 square and remove 1 from the first pile. Reset. Slide down 2 squares and remove 2 from the second pile. Choose two more slides; make each matching pile move.}
\Piles{1}
\begin{center}\renewcommand{\arraystretch}{1.8}\begin{tabular}{p{2.7in}|p{2.7in}}My slide from (3,2)&The two piles left\\\hline&\\&\\\end{tabular}\end{center}
\Problem{6}{Now play with piles only. On each turn remove one or more counters from one pile. Take the last counter to win. Try each start. If the piles differ, make them equal in one move. If they match, let your partner move and then restore the match. Record a move and reply.}
\begin{center}\renewcommand{\arraystretch}{2.25}\begin{tabular}{c|p{2.15in}|p{2.15in}}Start&After first move&After reply\\\hline(2,2)&&\\(3,2)&&\\(1,4)&&\\(4,4)&&\end{tabular}\end{center}
\newpage
\StudentPage{Grades 2--3}
\Problem{7}{Add a third pile; keep removing from only one pile per turn. Play (1,1,1), then (1,1,2). Draw a first move that leaves two matching piles and one empty pile in each game.}
\Space{1.1in}
\Problem{8}{Build (1,2,3). The table gives every possible first move. For each row, rebuild those piles and find a reply leaving two equal piles and one empty pile. Write the three piles after your reply. Keep the pile order unchanged.}
\begin{center}\renewcommand{\arraystretch}{2.4}\begin{tabular}{c|p{3.5in}}After the first move&After my reply\\\hline(0,2,3)&\\(1,0,3)&\\(1,1,3)&\\(1,2,0)&\\(1,2,1)&\\(1,2,2)&\\\end{tabular}\end{center}
\Problem{9}{Choose a two-pile start with unequal piles. Ask a partner to leave you equal piles in one move. Record your start and the move.}
\Space{.75in}

\end{document}
